\documentclass[11pt]{article}
\usepackage[margin=1in]{geometry}
\usepackage{hyperref}
\usepackage{listings}
\usepackage{xcolor}
\usepackage{amsmath}

\lstset{
  basicstyle=\ttfamily\small,
  breaklines=true,
  columns=fullflexible,
  frame=single,
  backgroundcolor=\color{gray!10},
  showstringspaces=false
}

\title{URI OPeNDAP Test Case: Working Log (Prompt \#7)}
\author{Compiled for P. Cornillon}
\date{2026-05-18}

\begin{document}
\maketitle

\section*{Purpose}

Working log for Prompt~\#7: write a proposal-grade summary of the
\texttt{grep-dap} test case and place it at the top of the project
repository, in a form that could feed a NASA Earth-science archive
metadata-validation proposal.

\section{Prompt \#7 (2026-05-18 21:57 UTC)}

\subsection{Instruction}

\emph{``Reread this document and documentation produced while
executing prompts 1--6 so that you are familiar with the work you
have already completed related to this project. Then write
\texttt{project\_summary.tex}, a \texttt{.tex} file summarising
what was done and why, and put the file in the top level of the
project. This file should be written in such a way that it could
be used as the basis for a portion of a proposal related to this
approach to help curators of NASA Earth Science archives to
validate and update the metadata of archives for which they are
responsible as well as to facilitate the use of these archives by
the research community.''}

\subsection{Drafting choices}

\paragraph{Location and naming.} Per the instruction the document
lives at the top level of the repository, not in
\texttt{docs/}: \texttt{/project\_summary.tex}.

\paragraph{Audience and tone.} The audience is a proposal
reviewer (or someone writing a section of a proposal). The tone
is therefore more declarative than the per-prompt reports: it
asserts what \texttt{grep-dap} demonstrates and what an extension
would do, rather than walking through the methodology a step at a
time. The detailed per-prompt write-ups in \texttt{docs/} carry
the methodology.

\paragraph{Structure.} The document follows a familiar
proposal-section progression:
\begin{enumerate}
  \item Abstract.
  \item Motivation: OPeNDAP's deliberate asymmetric metadata
    design, and the user / curator burden it creates.
  \item The grep-dap approach: three classes of evidence
    (catalog, structural, data probes); confidence-tagged
    inference; explicit refusal-to-guess.
  \item The URI case study (six subsections, one per prompt).
  \item Capabilities the case study demonstrates (six explicit
    capabilities).
  \item Implications for a NASA-scale service (catalog audit,
    cross-link inheritance, data-probe physical inference,
    dual-audience artefacts, evidence trail).
  \item Limitations and open questions (sample size, agent
    reliability, producer overhead, authentication, GenAI
    failure modes).
  \item Proposed next steps (multi-DAAC pilot,
    producer-loop integration, user-facing usage-doc service).
  \item Reproducibility (pointer to repository contents).
\end{enumerate}

\paragraph{Two errors from Prompt~1 are surfaced explicitly} in
the case-study section as a feature of the approach, not a
defect: a deployable service must catch and correct its own
prior inferences. Surfacing the corrections is a key part of the
trust argument for an automated audit service.

\paragraph{Claims I avoided making.}
\begin{itemize}
  \item No claim about NASA archive metadata distribution at
    scale; I only claim that the URI server documents a clear
    within-server bimodal distribution.
  \item No claim about specific NASA solicitations or programs.
  \item No claim about LLM model identity or performance
    benchmarks; the focus is on the methodology of how an agent
    is disciplined by data probes, not on which model is doing
    the inference.
\end{itemize}

\subsection{Deliverables}

\begin{itemize}
  \item \texttt{project\_summary.tex} --- the new top-level
    proposal summary.
  \item \texttt{docs/test\_case\_uri\_log\_7.tex} --- this log.
\end{itemize}

\subsection{Commands}

\begin{lstlisting}
pdflatex project_summary.tex
pdflatex test_case_uri_log_7.tex
\end{lstlisting}

\subsection{Session timing}

Started 2026-05-18 21:57\,UTC. Closed log 2026-05-18
$\sim$22:30\,UTC.

\end{document}
